\usepackage[T2A]{fontenc}
\usepackage[utf8]{inputenc}
\usepackage[english,russian]{babel}
\usepackage{tempora}
\usepackage{geometry}
\usepackage{setspace}
\usepackage{graphicx}
\usepackage[fleqn]{amsmath}
\usepackage{amssymb,mathtools}
\setlength{\mathindent}{1.25cm}
\usepackage{listings,xcolor}
\usepackage{hyperref}
\usepackage{caption,subcaption}
\usepackage{float}
\usepackage{indentfirst}
\usepackage{multirow,longtable,array,booktabs,tabularx}
\usepackage{makecell}
\usepackage{enumitem}
\usepackage{titlesec}
\usepackage{tocloft}

% --- Пути к рисункам проекта ---
\graphicspath{{project/figures/}{figures/}{./}}

% --- Нумерация страниц: внизу по центру ---
\pagestyle{plain}

% --- Оформление содержания ---
\renewcommand{\cfttoctitlefont}{\hfill\normalsize\bfseries}
\renewcommand{\cftaftertoctitle}{\hfill\mbox{}}
\setlength{\cftaftertoctitleskip}{1\baselineskip}
\renewcommand{\cftsecfont}{}
\renewcommand{\cftsecpagefont}{}
\renewcommand{\cftsubsecfont}{}
\renewcommand{\cftsubsecpagefont}{}
\renewcommand{\cftsubsubsecfont}{}
\renewcommand{\cftsubsubsecpagefont}{}
\renewcommand{\cftsecleader}{\cftdotfill{\cftdotsep}}
\renewcommand{\cftsubsecleader}{\cftdotfill{\cftdotsep}}
\renewcommand{\cftsubsubsecleader}{\cftdotfill{\cftdotsep}}
\setlength{\cftsecnumwidth}{1em}
\setlength{\cftsubsecnumwidth}{2.3em}
\setlength{\cftsubsubsecnumwidth}{2.5em}
\setlength{\cftsecindent}{0pt}
\setlength{\cftsubsecindent}{5mm}
\setlength{\cftsubsubsecindent}{10mm}
\setlength{\cftbeforesecskip}{0pt}
\setlength{\cftbeforesubsecskip}{0pt}
\setlength{\cftbeforesubsubsecskip}{0pt}
\cftsetpnumwidth{0.7em}
\cftsetrmarg{0.7em}
\setcounter{tocdepth}{2}

% --- Содержание: номер раздела + фиксированный пробел, без выравнивания
%     заголовков по левому краю (позиция названия зависит от ширины номера) ---
\makeatletter
\renewcommand{\numberline}[1]{#1\hspace{0.5em}}
\newcommand{\gost@tocline}[3]{%
  % #1 -- отступ слева, #2 -- запись (\numberline{..}Заголовок), #3 -- страница
  \vskip \z@ \@plus.2\p@
  {\leftskip #1\relax \rightskip \@tocrmarg \parfillskip -\rightskip
   \parindent #1\relax \@afterindenttrue
   \interlinepenalty\@M
   \leavevmode \normalsize
   #2\nobreak
   \leaders\hbox{$\m@th\mkern\@dotsep mu\hbox{.}\mkern\@dotsep mu$}\hfill
   \nobreak \hb@xt@\@pnumwidth{\hfil\normalcolor #3}\par}}
\renewcommand*\l@section{\gost@tocline{0pt}}
\renewcommand*\l@subsection{\gost@tocline{5mm}}
\renewcommand*\l@subsubsection{\gost@tocline{10mm}}
\makeatother

\geometry{left=30mm,right=10mm,top=20mm,bottom=20mm}
\onehalfspacing
\setlength{\parindent}{1.25cm}
\renewcommand{\arraystretch}{1.18}

% --- Заголовки нумерованных разделов: с отступом 1,25 см ---
\titleformat{\section}{\clearpage\normalsize\bfseries}{\thesection}{0.5em}{}
\titleformat{\subsection}{\normalsize\bfseries}{\thesubsection}{0.5em}{}
\titleformat{\subsubsection}{\normalsize\bfseries}{\thesubsubsection}{0.5em}{}
\titlespacing*{\section}{\parindent}{0pt}{\baselineskip}
\titlespacing*{\subsection}{\parindent}{\baselineskip}{\baselineskip}
\titlespacing*{\subsubsection}{\parindent}{\baselineskip}{\baselineskip}

% --- Заголовки ненумерованных разделов: по центру, с новой страницы ---
\titleformat{name=\section,numberless}{\clearpage\normalsize\bfseries\centering}{}{0em}{}
\titlespacing*{name=\section,numberless}{0pt}{0pt}{\baselineskip}

% --- Перечисления: отступ только на первой строке, тире вместо точек ---
\setlist{nolistsep}
\setlist[enumerate]{wide, itemindent=\parindent, listparindent=0pt}
\setlist[itemize]{wide, itemindent=\parindent, listparindent=0pt, label=-}

% --- Название библиографии ---
\renewcommand{\refname}{СПИСОК ИСПОЛЬЗОВАННЫХ ИСТОЧНИКОВ}

\DeclareCaptionLabelFormat{gost}{Рисунок #2}
\DeclareCaptionLabelFormat{gosttable}{Таблица #2}
\captionsetup[figure]{labelformat=gost,labelsep=endash,justification=centering}
\captionsetup[table]{labelformat=gosttable,labelsep=endash,justification=justified,singlelinecheck=false}

\definecolor{codegreen}{rgb}{0,0.6,0}
\definecolor{codegray}{rgb}{0.5,0.5,0.5}
\definecolor{codepurple}{rgb}{0.58,0,0.82}
\definecolor{backcolour}{rgb}{0.95,0.95,0.92}
\lstdefinestyle{pythonstyle}{
    backgroundcolor=\color{white},commentstyle=\color{codegreen},
    keywordstyle=\color{magenta},numberstyle=\tiny\color{codegray},
    stringstyle=\color{codepurple},basicstyle=\ttfamily\footnotesize,
    breaklines=true,captionpos=b,keepspaces=true,numbers=left,
    numbersep=5pt,showstringspaces=false,tabsize=2,language=Python,
    frame=single}
\lstset{style=pythonstyle}

\numberwithin{equation}{section}
\numberwithin{figure}{section}
\numberwithin{table}{section}

\hypersetup{colorlinks=true,linkcolor=black,filecolor=black,urlcolor=blue,citecolor=black}
